% Part: incompleteness
% Chapter: representability-in-q
% Section: composition-representable

\documentclass[../../../include/open-logic-section]{subfiles}

\begin{document}

\olfileid{inc}{req}{cmp}
\olsection{ترکیب در~$\Th{Q}$ بازنمایی‌پذیر است}

فرض کنید $h$ به‌صورت زیر تعریف شده باشد:
\[
h(x_0,\dots,x_{l-1}) = f(g_0(x_0,\dots,x_{l-1}), \dots,
g_{k-1}(x_0,\dots,x_{l-1})).
\]
و فرض کنید پیش‌تر !!{formula}sی $!A_f, !A_{g_0}, \dots,
!A_{g_{k-1}}$ را یافته‌ایم که به‌ترتیب توابع~$f$، $g_0$،
\dots،~$g_{k-1}$ را بازنمایی می‌کنند. باید !!a{formula}ی چون~$!A_h$
بیابیم که~$h$ را بازنمایی کند.

با حالتی ساده آغاز کنیم که همهٔ توابع $1$موضعی‌اند؛ یعنی
$h(x) = f(g(x))$ را در نظر بگیرید. اگر $!A_f(y, z)$ تابع~$f$ و
$!A_g(x, y)$ تابع~$g$ را بازنمایی کند، به !!a{formula}ی چون~$!A_h(x, z)$
نیاز داریم که~$h$ را بازنمایی کند. توجه کنید $h(x) = z$ اگر و تنها
اگر~$y$ای وجود داشته باشد که هم $z = f(y)$ و هم $y = g(x)$. (اگر
$h(x) = z$، آنگاه $g(x)$ چنین~$y$ای است؛ و اگر چنین~$y$ای وجود داشته
باشد، چون $y = g(x)$ و $z = f(y)$، داریم $z = f(g(x))$.) این نکته
نشان می‌دهد $\lexists[y][(!A_g(x, y) \land !A_f(y,
  z))]$ نامزد مناسبی برای~$!A_h(x, z)$ است. فقط باید بررسی کنیم که
$\Th{Q}$ !!{formula}sی مربوط را ثابت می‌کند.

\begin{prop}
\ollabel{prop:rep1}
اگر $h(n) = m$، آنگاه $\Th{Q} \Proves !A_h(\num{n}, \num{m})$.
\end{prop}

\begin{proof}
فرض کنید $h(n) = m$؛ یعنی $f(g(n)) = m$. بگذارید $k = g(n)$. آنگاه
\begin{align*}
  \Th{Q} & \Proves !A_g(\num{n}, \num{k})
  \intertext{زیرا $!A_g$ تابع~$g$ را بازنمایی می‌کند، و}
  \Th{Q} & \Proves !A_f(\num{k}, \num{m})
  \intertext{زیرا $!A_f$ تابع~$f$ را بازنمایی می‌کند. پس}
  \Th{Q} & \Proves !A_g(\num{n}, \num{k}) \land !A_f(\num{k}, \num{m})
  \intertext{و در نتیجه همچنین}
  \Th{Q} & \Proves \lexists[y][(!A_g(\num{n}, y) \land !A_f(y, \num{m}))],
\end{align*}
یعنی $\Th{Q} \Proves !A_h(\num{n}, \num{m})$.
\end{proof}

\begin{prop}
\ollabel{prop:rep2}
اگر $h(n) = m$، آنگاه $\Th{Q} \Proves \lforall[z][(!A_h(\num{n}, z) \lif
  z = \num{m})]$.
\end{prop}

\begin{proof}
فرض کنید $h(n) = m$؛ یعنی $f(g(n)) = m$. بگذارید $k = g(n)$. آنگاه
\begin{align*}
  \Th{Q} & \Proves \lforall[y][(!A_g(\num{n}, y) \lif \eq[y][\num{k}])]
  \intertext{زیرا $!A_g$ تابع~$g$ را بازنمایی می‌کند، و}
  \Th{Q} & \Proves \lforall[z][(!A_f(\num{k}, z) \lif \eq[z][\num{m}])]
  \intertext{زیرا $!A_f$ تابع~$f$ را بازنمایی می‌کند. تنها با اندکی
    منطق می‌توانیم نشان دهیم که همچنین}
  \Th{Q} & \Proves \lforall[z][(\lexists[y][(!A_g(\num{n}, y) \land
      !A_f(y, z))] \lif \eq[z][\num{m}])].
\end{align*}
یعنی $\Th{Q} \Proves \lforall[y][(!A_h(\num n, y) \lif \eq[y][\num m])]$.
\end{proof}

همین ایده در حالت پیچیده‌تری که $f$ و~$g_i$ موضع‌هایی بیش از~$1$
دارند نیز کار می‌کند.

\begin{prop}
\ollabel{prop:rep-composition}
اگر $!A_f(y_0, \dots, y_{k-1}, z)$ تابع $f(y_0, \dots, y_{k-1})$ را
در~$\Th{Q}$ بازنمایی کند، و $!A_{g_i}(x_0, \dots, x_{l-1}, y)$ تابع
$g_i(x_0, \dots, x_{l-1})$ را در~$\Th{Q}$ بازنمایی کند، آنگاه
\begin{multline*}
  \lexists[y_0\dots][\lexists[y_{k-1}][(!A_{g_0}(x_0,\dots,x_{l-1},y_0) \land
      \dots \land {}]]\\
  !A_{g_{k-1}}(x_0,\dots,x_{l-1},y_{k-1}) \land !A_f(y_0,\dots,y_{k-1},z))
\end{multline*}
تابع
\[
h(x_0, \dots, x_{l-1}) = f(g_0(x_0, \dots, x_{l-1}), \dots, g_{k-1}(x_0,
\dots, x_{l-1})).
\]
را بازنمایی می‌کند.
\end{prop}

\begin{proof}
تمرین.
\end{proof}

\begin{prob}
با راهنمایی‌گرفتن از اثبات‌های \olref[inc][req][cmp]{prop:rep1} و
\olref[inc][req][cmp]{prop:rep2}، اثبات
\olref[inc][req][cmp]{prop:rep-composition} را با جزئیات انجام دهید.
\end{prob}

\end{document}
